%% ============================================================
%%  DOCUMENT DE CADRAGE RÉSUMÉ — AGENT HYPOTHÈSES
%%  Version condensée 10 pages maximum
%%  Basé sur le document complet venv/cadrage.tex
%% ============================================================
\documentclass[11pt, a4paper]{article}

%% ---------- Encodage & Langue ----------
\usepackage[utf8]{inputenc}
\usepackage[T1]{fontenc}
\usepackage[french]{babel}

%% ---------- Mise en page compacte ----------
\usepackage[top=2cm, bottom=2cm, left=2cm, right=2cm]{geometry}
\usepackage{parskip}
\setlength{\parindent}{0pt}
\setlength{\parskip}{4pt}

%% ---------- Packages ----------
\usepackage{booktabs}
\usepackage{tabularx}
\usepackage{xcolor}
\usepackage{tcolorbox}
\usepackage{amsmath}
\usepackage{amssymb}
\usepackage{graphicx}
\usepackage{hyperref}
\usepackage{fancyvrb}

%% ---------- Couleurs ----------
\definecolor{navyblue}{HTML}{1A3A5C}
\definecolor{midblue}{HTML}{2E6DA4}
\definecolor{lightblue}{HTML}{D6E8F7}
\definecolor{paleblue}{HTML}{EBF4FB}
\definecolor{darkgreen}{HTML}{1B5E3B}
\definecolor{lightgreen}{HTML}{D4EDDA}

%% ---------- Boîtes colorées ----------
\newtcolorbox{formulebox}[1][]{
  colback=paleblue, colframe=midblue, 
  fonttitle=\bfseries\color{navyblue},
  boxrule=0.8pt, arc=3pt,
  left=6pt, right=6pt, top=4pt, bottom=4pt,
  #1
}

\newtcolorbox{resultbox}[1][]{
  colback=lightgreen, colframe=darkgreen,
  boxrule=0.6pt, arc=3pt,
  left=6pt, right=6pt, top=4pt, bottom=4pt,
  #1
}

%% ---------- Formatage titres ----------
\usepackage{titlesec}
\titleformat{\section}[block]
  {\normalfont\Large\bfseries\color{navyblue}}
  {\thesection}{0.5em}{}
  [\vspace{2pt}\textcolor{midblue}{\hrule height 0.8pt}]

\titleformat{\subsection}[block]
  {\normalfont\large\bfseries\color{midblue}}
  {\thesubsection}{0.4em}{}

%% ============================================================
\begin{document}

%% ============================================================
%%  PAGE DE GARDE
%% ============================================================
\begin{titlepage}
\pagecolor{navyblue}
\color{white}
\begin{center}

\vspace*{1.5cm}

{\small\bfseries\color{lightblue} UNIVERSITÉ ABDELMALEK ESSAÂDI — FSTT}\\[4pt]
{\small\color{lightblue} Licence Data Analytics}\\[6pt]
{\small\color{lightblue} Projet de Fin d'Études}

\vspace{1.2cm}
\textcolor{midblue}{\hrule height 1pt}
\vspace{0.3cm}
\textcolor{rulegray}{\hrule height 0.5pt}
\vspace{1.2cm}

{\fontsize{13}{18}\selectfont\bfseries\color{lightblue}
DOCUMENT DE CADRAGE LOGIQUE}\\[10pt]

{\fontsize{26}{34}\selectfont\bfseries\color{white}
Agent Hypothèses}\\[14pt]

{\fontsize{14}{20}\selectfont\color{lightblue}
Module Pédagogique Intelligent}\\[6pt]

{\fontsize{11}{16}\selectfont\color{rulegray}
Plateforme IA End-to-End :
Idéation $\rightarrow$ Business Plan $\rightarrow$ Immatriculation}

\vspace{1.2cm}
\textcolor{rulegray}{\hrule height 0.5pt}
\vspace{0.3cm}
\textcolor{midblue}{\hrule height 1pt}
\vspace{1.8cm}

%% Trois blocs thématiques
\begin{minipage}{0.3\textwidth}
\centering
\textcolor{midblue}{\large\bfseries Chapitre 1}\\[4pt]
\textcolor{lightblue}{\small Introduction \&\\Contexte Stratégique}
\end{minipage}
\hfill
\begin{minipage}{0.3\textwidth}
\centering
\textcolor{midblue}{\large\bfseries Chapitre 2}\\[4pt]
\textcolor{lightblue}{\small Notions Fondamentales\\\& Logique Métier}
\end{minipage}
\hfill
\begin{minipage}{0.3\textwidth}
\centering
\textcolor{midblue}{\large\bfseries Chapitre 3}\\[4pt]
\textcolor{lightblue}{\small Veille Technologique\\\& Stack Technique}
\end{minipage}

\vfill

{\small\color{rulegray} Année Universitaire 2025--2026}

\end{center}
\end{titlepage}

\nopagecolor
\color{black}

%% ============================================================
%%  TABLE DES MATIÈRES
%% ============================================================
\tableofcontents
\newpage

\nopagecolor
\color{black}

%% ============================================================
%%  CHAPITRE 1 — INTRODUCTION & CONTEXTE
%% ============================================================
\section{Introduction \& Contexte Stratégique}

\subsection{Présentation du projet}

L'Agent Hypothèses constitue le cœur de la plateforme intelligente d'accompagnement à la création d'entreprise au Maroc. Développé dans le cadre d'un Projet de Fin d'Études en Licence Data Analytics (FSTT), ce module IA orchestre un dialogue pédagogique avec l'entrepreneur pour collecter \textbf{22 hypothèses financières} structurées selon la méthode Chauvin, les valider automatiquement et générer 3 scénarios financiers prévisionnels.

\subsection{Position dans l'écosystème multi-agents}

L'Agent Hypothèses s'inscrit dans la \textbf{Phase 2 — Stratégie \& Business Plan} du workflow global :

\begin{formulebox}[title={\bfseries\color{navyblue} Workflow complet}]
\texttt{Angular $\rightarrow$ FastAPI $\rightarrow$ SLM Phi-3 Mini}\\
\texttt{$\rightarrow$ Agent Contrôleur n8n $\rightarrow$ POST /hypotheses/start}\\
\texttt{$\rightarrow$ Dialogue 22 questions (RAG + Groq)}\\
\texttt{$\rightarrow$ JSON $\rightarrow$ Agent Finance $\rightarrow$ PDF final}
\end{formulebox}

\subsection{Caractéristiques techniques principales}

\begin{center}
\begin{tabularx}{0.9\textwidth}{X X}
\toprule
\rowcolor{navyblue}\textcolor{white}{\textbf{Composant}} & \textcolor{white}{\textbf{Technologie}} \\
\midrule
LLM & Groq API (llama-3.3-70b-versatile) \\
Base vectorielle & ChromaDB (280 documents) \\
Embeddings & sentence-transformers (all-MiniLM-L6-v2) \\
API REST & FastAPI (port 8000) \\
Validation & Règles Chauvin codées en Python \\
\bottomrule
\end{tabularx}
\end{center}

%% ============================================================
%%  CHAPITRE 2 — NOTIONS FONDAMENTALES
%% ============================================================
\section{Notions Fondamentales \& Logique Métier}

\subsection{Architecture RAG (Retrieval-Augmented Generation)}

Le pipeline RAG combine recherche sémantique et génération de texte :

\begin{enumerate}
  \item \textbf{Embedding} : Conversion du texte en vecteur (384 dimensions)
  \item \textbf{Recherche} : Similarité cosinus dans ChromaDB (280 documents)
  \item \textbf{Génération} : Enrichissement du prompt LLM avec contexte pertinent
\end{enumerate}

\subsection{Les 22 hypothèses en 4 blocs}

\begin{center}
\begin{tabularx}{0.95\textwidth}{X X X}
\toprule
\rowcolor{navyblue}\textcolor{white}{\textbf{Bloc}} & \textcolor{white}{\textbf{Question Chauvin}} & \textcolor{white}{\textbf{Hypothèses}} \\
\midrule
1 & Quoi ? & Ventes, clients, prix, croissance (H1-H7) \\
2 & Comment ? & Type activité, fabrication, fournisseurs (H8-H12) \\
3 & Combien ? & Charges fixes, loyer, salaires, marketing (H13-H21) \\
4 & Quand ? & Nature clients, délais encaissement (H22) \\
\bottomrule
\end{tabularx}
\end{center}

\subsection{Règles de validation de Chauvin}

\begin{formulebox}[title={\bfseries\color{navyblue} Règle 1 — Seuil de rentabilité}]
\[
\text{Seuil clients} = \frac{\text{Charges fixes}}{\text{Prix} - \text{Coût fabrication}}
\]
\textbf{Alerte} si clients prévus $<$ seuil calculé
\end{formulebox}

\begin{formulebox}[title={\bfseries\color{navyblue} Règle 2 — Symétrie financement}]
\textbf{Alerte} si investissements $> 2 \times$ apports propres
\end{formulebox}

\begin{formulebox}[title={\bfseries\color{navyblue} Règle 3 — Cohérence prix/coût}]
\textbf{Alerte} si coût fabrication $> 80\%$ prix de vente
\end{formulebox}

\subsection{Génération des 3 scénarios}

Application automatique des coefficients :
\begin{itemize}
  \item \textbf{Pessimiste} : $\times 0.7$ (30\% moins de clients)
  \item \textbf{Réaliste} : $\times 1.0$ (hypothèses saisies)
  \item \textbf{Optimiste} : $\times 1.3$ (30\% plus de clients)
\end{itemize}

Pour chaque scénario : CA mensuel 12 mois, charges, marge, résultat net, mois de rentabilité.

%% ============================================================
%%  CHAPITRE 3 — VEILLE TECHNOLOGIQUE
%% ============================================================
\section{Veille Technologique \& Stack Technique}

\subsection{Choix technologiques justifiés}

\begin{center}
\begin{tabularx}{0.95\textwidth}{X X X}
\toprule
\rowcolor{navyblue}\textcolor{white}{\textbf{Catégorie}} & \textcolor{white}{\textbf{Choix}} & \textcolor{white}{\textbf{Justification}} \\
\midrule
LLM & Groq API & Ultra-rapide (LPU), gratuit, llama-3.3-70b \\
Base vectorielle & ChromaDB & Open-source, Python natif, gratuit \\
Embeddings & all-MiniLM-L6-v2 & Léger (80 Mo), performant, 384 dims \\
API & FastAPI & Performance, documentation auto, Pydantic \\
Orchestration & n8n & Workflow multi-agents, open-source \\
\bottomrule
\end{tabularx}
\end{center}

\subsection{Architecture modulaire}

\begin{resultbox}
\textbf{4 couches principales :}\\
\textbf{API} (\texttt{main.py}) : Point d'entrée HTTP, 5 endpoints\\
\textbf{Agent} (\texttt{agent/}) : Dialogue, validation, construction JSON\\
\textbf{RAG} (\texttt{rag/}) : Recherche sémantique, 280 documents\\
\textbf{Schémas} (\texttt{schemas/}) : Contrat JSON Pydantic
\end{resultbox}

\subsection{Arborescence du projet}

\begin{Verbatim}[fontsize=\small,frame=single]
agent_hypotheses/
├── .env                    # Variables d'environnement
├── main.py                 # API FastAPI principale
├── requirements.txt        # Dépendances Python
├── agent/                  # Logique métier
│   ├── dialogue.py        # Cœur conversationnel
│   ├── groq_client.py     # Connexion LLM
│   ├── validation.py      # Règles de Chauvin
│   ├── questions.py       # 22 questions structurées
│   ├── scenarios.py       # 3 projections financières
│   └── output_builder.py  # Construction JSON final
├── rag/                    # Pipeline RAG
│   ├── vectorstore.py      # ChromaDB
│   ├── embedder.py         # Modèle d'embeddings
│   └── retriever.py        # Recherche sémantique
├── schemas/                # Schémas Pydantic
│   └── output_schema.py    # Contrat JSON
├── data/                   # Données de référence
│   └── scenarios_reference.json
└── scripts/                # Tests et validation
    ├── test_dialogue.py
    ├── test_validation.py
    ├── test_json_final.py
    └── test_scenarios_rag.py
\end{Verbatim}

\subsection{Métriques de performance}

\begin{center}
\begin{tabularx}{0.9\textwidth}{X X}
\toprule
\rowcolor{navyblue}\textcolor{white}{\textbf{Indicateur}} & \textcolor{white}{\textbf{Résultat}} \\
\midrule
Connexion Groq & 1 070 ms ($<$ 2s requis) \\
Documents RAG & 280 indexés (270 Chauvin + 10 scénarios) \\
Questions couvertes & 22 / 22 \\
Endpoints API & 5 / 5 fonctionnels \\
Scénarios générés & 3 automatiques \\
JSON final & 12 champs validés \\
\bottomrule
\end{tabularx}
\end{center}

\subsection{Exemple concret — Épicerie Tanger}

\textbf{Hypothèses} : Prix 150 MAD, 30 clients, coût 80 MAD, loyer 3 000 MAD, salaires 6 000 MAD

\begin{center}
\begin{tabularx}{0.95\textwidth}{X X X X}
\toprule
\rowcolor{navyblue}\textcolor{white}{\textbf{Indicateur}} & \textcolor{white}{\textbf{Pessimiste}} & \textcolor{white}{\textbf{Réaliste}} & \textcolor{white}{\textbf{Optimiste}} \\
\midrule
CA annuel & 67 360 MAD & 96 229 MAD & 125 098 MAD \\
Résultat net an 1 & $-$112 465 MAD & $-$98 993 MAD & $-$85 521 MAD \\
Mois rentabilité & Mois 22 & Mois 19 & Mois 16 \\
\bottomrule
\end{tabularx}
\end{center}

\begin{resultbox}
\textbf{Message pédagogique généré :} Votre projet nécessite 19 mois avant rentabilité dans le scénario réaliste. Assurez-vous d'avoir une trésorerie suffisante pour couvrir les 19 premiers mois d'activité.
\end{resultbox}

%% ============================================================
%%  CONCLUSION
%% ============================================================
\section{Conclusion}

L'Agent Hypothèses constitue une solution innovante qui transforme l'accompagnement entrepreneurial au Maroc en combinant intelligence artificielle, méthodologie financière rigoureuse et pédagogie adaptée. En collectant structurément les 22 hypothèses financières selon la méthode Chauvin, en les validant automatiquement et en générant trois scénarios prévisionnels, ce module IA démocratise l'accès à des outils financiers de qualité pour les entrepreneurs marocains. L'intégration d'un pipeline RAG avec 280 documents de référence et d'une API REST performante garantit la fiabilité et l'accessibilité de la solution, positionnant l'Agent Hypothèses comme un pont essentiel entre l'idée entrepreneuriale et la réalisation d'un business plan solide et validé.

\end{document}
